\documentclass{article}
\usepackage{graphicx} % Required for inserting images
\usepackage{amsmath}
\usepackage{amssymb}
\usepackage{amsthm}
%\usepackage{hyperref}

\title{Finite Complement Topology}
\author{Aaron Honculada}
\date{February 2025}

\begin{document}

\maketitle

\section{Definition}

The \textbf{finite complement topology} on a set $X$ is defined as follows:

\begin{itemize}
    \item Open sets are:
    \begin{itemize}
        \item The empty set $\emptyset$.
        \item Any subset of $X$ whose complement is finite(i.e., contains all but finitely many elements of $X$).
    \end{itemize}
    \item Closed sets are:
    \begin{itemize}
        \item The whole space $X$.
        \item Any finite subset of $X$.
    \end{itemize}
\end{itemize}

\section{Example: Finite Complement Topology on $\mathbb{R}$}
Let $X=\mathbb{R}$ with the finite complement topology. Then:

\begin{itemize}
    \item The set $(-\infty, 0)$ is not open because its complement $[0, \infty)$ is not finite.
    \item The set $\mathbb{R}\setminus \{1,2,3\}$ is open because its complement $\{1,2,3\}$ is finite.
    \item Any finite set, like $\{5,7,9\}$, is \textbf{closed}.
\end{itemize}

\section{Example: Finite Complement Topology on $\mathbb{Z}$}
Let $X=\mathbb{Z}$ be the set of integers, and define a topology $\tau$ where:  

\begin{itemize}
    \item The open sets in $\mathbb{Z}$ are:
    \begin{itemize}
        \item The empty set $\emptyset$.
        \item The whole space $\mathbb{Z}$.
        \item Any subset of $\mathbb{Z}$ whose complement is finite, meaning a set $U \subseteq
        \mathbb{Z}$ is open if and only if $\mathbb{Z} \setminus U$ is finite.
    \end{itemize}
    \item The closed sets in $\mathbb{Z}$ are:
    \begin{itemize}
        \item The whole space $\mathbb{Z}$.
        \item Any finite subset of $\mathbb{Z}$.
    \end{itemize}
\end{itemize}

\begin{itemize}
    \item The set $\{0,1,2\}$ is not open because its complement $\mathbb{Z}\setminus \{0,1,2\}$ is not finite.
\end{itemize}
\subsection{Why is this interesting?}
\begin{itemize}
    \item The space $(\mathbb{Z}, \tau)$ is \textbf{not Hausdorff} because
    distinct points cannot be separated by disjoint open sets.
    \item For example, the points $0$ and $1$ cannot be separated by disjoint open sets because
    the intersection of any open set containing $0$ and any open set containing $1$ is non-empty
    (it contains $0$ and $1$).
\end{itemize}

\section{Example: Finite Complement Topology on a finite set}
Let $X=\{a,b,c,d\}$ be a finite set, and define a topology $\tau$ where:

\begin{itemize}
    \item The open sets in $X$ are:
    \begin{itemize}
        \item The empty set $\emptyset$.
        \item The whole space $X=\{a,b,c,d\}$.
        \item Any subset $U$ such that $X\setminus U$ is \textbf{finite}.
    \end{itemize}
    \item Since $X$ is already finite:
    The complement of any subset of $X$ is always finite!
    This means that \textbf{every subset} of $X$ is open, so:
    \begin{itemize}
        \item $\tau = \mathcal{P}(X)= \text{the discrete topology}$
    \end{itemize}
\end{itemize}

\section{Example: Finite Complement Topology on $\mathbb{R}$}
Let $X=\mathbb{R}$ with the finite complement topology. Then:


\end{document}










